\par\Needspace{6\baselineskip}\medskip
\hypertarget{native-128-article095}{}\pdfbookmark[2]{95.}{article-095}
\begin{center}95.\end{center}
Transimus ad negotium secundum, puta determinationem elementorum ex invento angulo $g$. Semiaxis maior hic statim habetur per formulas 10, $10^{\ast}$, pro quibus etiam sequentes adhiberi possunt:
\[
\begin{aligned}
\mathrm{17)}\qquad a&=\dfrac{2mm\cos f\mathpunct{.}\BareRoot{rr'}}{yy\sin g^2}=\dfrac{kktt}{4yyrr'\cos f^2\sin g^2}\\[6pt]
\mathrm{17}^{\ast}\mathrm{)}\qquad a&=\dfrac{-2MM\cos f\mathpunct{.}\BareRoot{rr'}}{YY\sin g^2}=\dfrac{kktt}{4YYrr'\cos f^2\sin g^2}\mathpunct{\cdot}
\end{aligned}
\]
Semiaxis minor $b=\BareRoot{ap}$ habetur per aequationem 1, qua cum praecedentibus combinata prodit
\[
\begin{aligned}
\mathrm{18)}\qquad p&=\left(\dfrac{yrr'\sin 2f}{kt}\right)^2\\[6pt]
\mathrm{18}^{\ast}\mathrm{)}\qquad p&=\left(\dfrac{Yrr'\sin 2f}{kt}\right)^2\mathpunct{\cdot}
\end{aligned}
\]
Iam sector ellipticus inter duos radios vectores atque arcum ellipticum contentus fit $=\frac12 kt\BareRoot{p}$, triangulum autem inter eosdem radios vectores atque chordam $=\frac12 rr'\sin 2f$: quamobrem ratio sectoris ad triangulum est ut $y:1$ vel $Y:1$. Haec observatio maximi est momenti, simulque aequationes 12, $12^{\ast}$ pulcherrime illustrat: patet enim hinc, in aequatione 12 partes $m$, $(l+x)^{\frac12}$, $X(l+x)^{\frac32}$, in aequatione $12^{\ast}$ autem partes $M$, $(L-x)^{\frac12}$, $X(L-x)^{\frac32}$ respective proportionales esse areae sectoris (inter radios vectores atque arcum ellipticum), areae trianguli (inter radios vectores atque chordam), areae segmenti (inter arcum atque chordam), quoniam manifesto area prima aequalis est vel summae vel differentiae duarum reliquarum, prout $v'-v$ vel inter 0 et $180^0$ iacet vel inter $180^0$ et $360^0$. In casu eo, ubi $v'-v$ maior est quam $360^0$, areae \SourceLeaf{129}sectoris nec non areae segmenti aream integrae ellipsis toties adiectam concipere oportet, quot revolutiones integras ille motus continet.

Quum $b$ sit $=a\cos\varphi$, e combinatione aequationum 1, 10, $10^{\ast}$ porro sequitur
\[
\begin{aligned}
\mathrm{19)}\qquad \cos\varphi&=\dfrac{\sin g\operatorname{tang}f}{2(l+\sin\frac12 g^2)}\\[6pt]
\mathrm{19}^{\ast}\mathrm{)}\qquad \cos\varphi&=\dfrac{-\sin g\operatorname{tang}f}{2(L-\sin\frac12 g^2)}\mathpunct{\raisebox{.45ex}{,}}
\end{aligned}
\]
unde substituendo pro $l$, $L$ valores suos ex art. 89 prodit
\[
\mathrm{20)}\qquad \cos\varphi=\dfrac{\sin f\sin g}{1-\cos f\cos g+2\operatorname{tang}2\omega^2}\mathpunct{\cdot}
\]
Haec formula ad calculum exactum excentricitatis non est idonea, quoties haecce modica est: sed facile ex ista deducitur formula aptior sequens
\[
\mathrm{21)}\qquad \operatorname{tang}\tfrac12\varphi^2=\dfrac{\sin\frac12(f-g)^2+\operatorname{tang}2\omega^2}{\sin\frac12(f+g)^2+\operatorname{tang}2\omega^2}\mathpunct{\raisebox{.45ex}{,}}
\]
cui etiam forma sequens tribui potest (multiplicando numeratorem et denominatorem per $\cos 2\omega^2$)
\[
\mathrm{22)}\qquad \operatorname{tang}\tfrac12\varphi^2=\dfrac{\sin\frac12(f-g)^2+\cos\frac12(f-g)^2\sin 2\omega^2}{\sin\frac12(f+g)^2+\cos\frac12(f+g)^2\sin 2\omega^2}\mathpunct{\cdot}
\]
Per utramque formulam (adhibitis si placet angulis auxiliaribus quorum tangentes $\dfrac{\operatorname{tang}2\omega}{\sin\frac12(f-g)}\mathpunct{\raisebox{.45ex}{,}}$ $\dfrac{\operatorname{tang}2\omega}{\sin\frac12(f+g)}$ pro priori, vel $\dfrac{\sin 2\omega}{\operatorname{tang}\frac12(f-g)}\mathpunct{\raisebox{.45ex}{,}}$ $\dfrac{\sin 2\omega}{\operatorname{tang}\frac12(f+g)}$ pro posteriori) angulum $\varphi$ omni semper praecisione determinare licebit.

Pro determinatione anguli $G$ adhiberi potest formula sequens, quae sponte demanat e combinatione aequationum 5, 7 et sequentis non numeratae:
\[
\mathrm{23)}\qquad \operatorname{tang}G=\dfrac{(r'-r)\sin g}{(r'+r)\cos g-2\cos f\mathpunct{.}\BareRoot{rr'}}\mathpunct{\raisebox{.45ex}{,}}
\]
e qua, introducendo $\omega$, facile derivatur
\[
\mathrm{24)}\qquad \operatorname{tang}G=\dfrac{\sin g\sin 2\omega}{\cos 2\omega^2\sin\frac12(f-g)\sin\frac12(f+g)+\sin 2\omega^2\cos g}\mathpunct{\cdot}
\]
Ambiguitas hic remanens facile deciditur adiumento aequationis 7, quae docet, $G$ inter 0 et $180^0$ vel inter $180^0$ et $360^0$ accipi debere, prout numerator in his duabus formulis positivus fuerit vel negativus.

\SourceLeaf{130}Combinando aequationem 3 cum his, quae protinus demanant ex aequatione II art. 8,
\[
\begin{aligned}
\dfrac1r-\dfrac1{r'}&=\dfrac{2e}{p}\sin f\sin F\\[6pt]
\dfrac1r+\dfrac1{r'}&=\dfrac2p+\dfrac{2e}{p}\cos f\cos F,
\end{aligned}
\]
nullo negotio derivabitur sequens:
\[
\mathrm{25)}\qquad \operatorname{tang}F=\dfrac{(r'-r)\sin f}{2\cos g\mathpunct{.}\BareRoot{rr'}-(r'+r)\cos f}\mathpunct{\raisebox{.45ex}{,}}
\]
e qua, introducto angulo $\grave{\omega}$, prodit
\[
\mathrm{26)}\qquad \operatorname{tang}F=\dfrac{\sin f\sin 2\omega}{\cos 2\omega^2\sin\frac12(f-g)\sin\frac12(f+g)-\sin 2\omega^2\cos f}\mathpunct{\cdot}
\]
Ambiguitas hic perinde tollitur ut ante. — Postquam anguli $F$ et $G$ inventi erunt, habebitur $v=F-f$, $v'=F+f$, unde positio perihelii nota \raisebox{.45ex}{.}erit; nec non $E=G-g$, $E'=G+g$. Denique motus medius intra tempus $t$ erit $=\dfrac{kt}{a^{\frac32}}=2g-2e\cos G\sin g$, quarum expressionum consensus calculo confirmando inserviet; epocha autem anomaliae mediae, respondens temporis momento inter duo proposita medio, erit $G-e\sin G\cos g$, quae pro lubitu ad quodvis aliud tempus transferri poterit. Aliquanto adhuc commodius est, anomalias medias pro duobus temporum momentis datis per formulas $E-e\sin E$, $E'-e\sin E'$ computare, harumque differentia cum $\dfrac{kt}{a^{\frac32}}$ comparanda ad calculi confirmationem uti.
